\chapter{高等数学：三重积分}

\section{三重积分的极限、区域与坐标变换}

\begin{problemblock}
\textbf{例1.} 求极限
\[
\lim_{n\to\infty}\frac{\pi}{2n^5}
\sum_{i=1}^{n}\sum_{j=1}^{n}\sum_{k=1}^{n}
i^2\sin\frac{\pi j}{2n}\cos\frac{k}{n}.
\]
\end{problemblock}
\begin{solutionblock}
\analysis{三重求和可以拆成三个一维求和：
\[
\frac{\pi}{2n^5}
\left(\sum_{i=1}^n i^2\right)
\left(\sum_{j=1}^n\sin\frac{\pi j}{2n}\right)
\left(\sum_{k=1}^n\cos\frac{k}{n}\right).
\]
整理成 Riemann 和：
\[
\frac{\pi}{2}
\left(\frac1{n^3}\sum_{i=1}^n i^2\right)
\left(\frac1n\sum_{j=1}^n\sin\frac{\pi j}{2n}\right)
\left(\frac1n\sum_{k=1}^n\cos\frac{k}{n}\right).
\]
三个极限分别为
\[
\frac1{n^3}\sum_{i=1}^n i^2\to\int_0^1x^2\,dx=\frac13,
\]
\[
\frac1n\sum_{j=1}^n\sin\frac{\pi j}{2n}\to
\int_0^1\sin\frac{\pi y}{2}\,dy=\frac2\pi,
\]
\[
\frac1n\sum_{k=1}^n\cos\frac{k}{n}\to
\int_0^1\cos z\,dz=\sin1.
\]
故原极限为
\[
\frac{\pi}{2}\cdot\frac13\cdot\frac2\pi\cdot\sin1
=\frac{\sin1}{3}.
\]}
\answer{\(\displaystyle \frac{\sin1}{3}\)}
\examnote{多重求和若被乘积型通项分离，先拆和，再分别转化为一维 Riemann 和。}
\end{solutionblock}

\begin{problemblock}
\textbf{例2.} 确定 \(\Omega\) 的形状，使得三重积分
\[
I=\iiint_{\Omega}(1-x^2-2y^2-3z^2)\,dv
\]
取得最大值。
\end{problemblock}
\begin{solutionblock}
\analysis{如果积分区域 \(\Omega\) 可以自由选取，为使
\[
\iiint_{\Omega}f(x,y,z)\,dv
\]
最大，应取所有使 \(f(x,y,z)>0\) 的点，舍去使 \(f(x,y,z)<0\) 的点。这里
\[
f(x,y,z)=1-x^2-2y^2-3z^2.
\]
因此最优区域为
\[
\Omega=\{(x,y,z):1-x^2-2y^2-3z^2\ge0\},
\]
即椭球
\[
x^2+2y^2+3z^2\le1.
\]
若顺便求最大值，令
\[
u=x,\qquad v=\sqrt2\,y,\qquad w=\sqrt3\,z,
\]
则 \(dx\,dy\,dz=\frac1{\sqrt6}\,du\,dv\,dw\)，区域化为单位球 \(u^2+v^2+w^2\le1\)。最大值为
\[
\frac1{\sqrt6}\iiint_{u^2+v^2+w^2\le1}(1-u^2-v^2-w^2)\,du\,dv\,dw
\]
\[
=\frac1{\sqrt6}\cdot4\pi\int_0^1(1-r^2)r^2\,dr
=\frac{8\pi}{15\sqrt6}.
\]}
\answer{\(\Omega\) 应取椭球 \(x^2+2y^2+3z^2\le1\)。此时最大值为 \(\displaystyle \frac{8\pi}{15\sqrt6}\)。}
\examnote{“自由选区域使积分最大”本质是正负区域取舍：正的全取，负的全舍，零边界不影响积分值。}
\end{solutionblock}

\begin{problemblock}
\textbf{例3.} 计算
\[
I=\iiint_{\Omega}(x+y+z)\,dx\,dy\,dz,
\]
其中 \(\Omega\) 为三个坐标面及平面 \(x+y+z=1\) 所围成的区域。
\end{problemblock}
\begin{solutionblock}
\analysis{区域为第一卦限中的四面体
\[
\Omega=\{(x,y,z):x\ge0,\ y\ge0,\ z\ge0,\ x+y+z\le1\}.
\]
直接积分：
\[
I=\int_0^1\int_0^{1-x}\int_0^{1-x-y}(x+y+z)\,dz\,dy\,dx.
\]
也可用对称性。该四面体体积为
\[
V=\frac16.
\]
四面体的重心为
\[
\left(\frac14,\frac14,\frac14\right).
\]
因此
\[
\iiint_{\Omega}x\,dV=\frac16\cdot\frac14=\frac1{24},
\]
同理 \(y,z\) 的积分也各为 \(\frac1{24}\)。故
\[
I=3\cdot\frac1{24}=\frac18.
\]}
\method{方法二：变量代换}
令
\[
u=x+y+z.
\]
在标准单纯形中，截面 \(x+y+z=u\) 的面积随 \(u^2\) 成比例，也可得到
\[
I=\int_0^1u\cdot\frac{u^2}{2}\,du=\frac18.
\]
\answer{\(\displaystyle \frac18\)}
\examnote{标准四面体 \(x+y+z\le1\) 的体积是 \(1/6\)，重心坐标是 \((1/4,1/4,1/4)\)，这两个结论在考研中很省时间。}
\end{solutionblock}

\begin{problemblock}
\textbf{例4.} 计算
\[
\iiint_{\Omega}z\,dx\,dy\,dz,
\]
其中 \(\Omega\) 由
\[
z=\sqrt{2-x^2-y^2},\qquad z=x^2+y^2
\]
围成。
\end{problemblock}
\begin{solutionblock}
\analysis{两个曲面分别是上半球
\[
x^2+y^2+z^2=2,\qquad z\ge0,
\]
与抛物面
\[
z=x^2+y^2.
\]
用柱坐标 \(x=r\cos\theta,\ y=r\sin\theta\)。交线满足
\[
z=r^2,\qquad z=\sqrt{2-r^2}.
\]
于是
\[
r^2=\sqrt{2-r^2}
\quad\Longrightarrow\quad r^4+r^2-2=0,
\]
得 \(r^2=1\)，即 \(0\le r\le1\)。区域为
\[
0\le\theta\le2\pi,\qquad
0\le r\le1,\qquad
r^2\le z\le\sqrt{2-r^2}.
\]
因此
\[
I=\int_0^{2\pi}\int_0^1\int_{r^2}^{\sqrt{2-r^2}}z\,r\,dz\,dr\,d\theta.
\]
先对 \(z\) 积分：
\[
I=2\pi\int_0^1\frac12\left[(2-r^2)-r^4\right]r\,dr
\]
\[
=\pi\int_0^1(2r-r^3-r^5)\,dr
=\pi\left(1-\frac14-\frac16\right)
=\frac{7\pi}{12}.
\]}
\answer{\(\displaystyle \frac{7\pi}{12}\)}
\examnote{球面与旋转抛物面围成的区域，柱坐标通常最直接；交线先用 \(r\) 解出来。}
\end{solutionblock}

\begin{problemblock}
\textbf{例5.} 求
\[
\iiint_{\Omega}\sqrt{x^2+y^2+z^2}\,dV,
\]
其中 \(\Omega\) 由
\[
z=\sqrt{x^2+y^2}
\]
与平面 \(z=1\) 围成。
\end{problemblock}
\begin{solutionblock}
\analysis{曲面 \(z=\sqrt{x^2+y^2}\) 是圆锥 \(z=r\)，平面 \(z=1\) 截出锥体。由于被积函数是到原点的距离，用球坐标最方便：
\[
x=\rho\sin\varphi\cos\theta,\quad
y=\rho\sin\varphi\sin\theta,\quad
z=\rho\cos\varphi.
\]
圆锥 \(z=r\) 对应
\[
\varphi=\frac{\pi}{4},
\]
平面 \(z=1\) 对应
\[
\rho\cos\varphi=1,\qquad \rho=\sec\varphi.
\]
区域为
\[
0\le\theta\le2\pi,\qquad
0\le\varphi\le\frac{\pi}{4},\qquad
0\le\rho\le\sec\varphi.
\]
被积函数为 \(\rho\)，体积元为 \(\rho^2\sin\varphi\,d\rho\,d\varphi\,d\theta\)。故
\[
I=\int_0^{2\pi}\int_0^{\pi/4}\int_0^{\sec\varphi}
\rho^3\sin\varphi\,d\rho\,d\varphi\,d\theta
\]
\[
=\frac{\pi}{2}\int_0^{\pi/4}\sec^4\varphi\sin\varphi\,d\varphi.
\]
令 \(u=\cos\varphi\)，则
\[
\int_0^{\pi/4}\sec^4\varphi\sin\varphi\,d\varphi
=\int_{\sqrt2/2}^{1}u^{-4}\,du
=\frac{2\sqrt2-1}{3}.
\]
因此
\[
I=\frac{\pi(2\sqrt2-1)}6.
\]}
\answer{\(\displaystyle \frac{\pi(2\sqrt2-1)}6\)}
\examnote{被积函数含 \(\sqrt{x^2+y^2+z^2}\) 时，优先考虑球坐标；圆锥给 \(\varphi\) 的范围，水平面给 \(\rho\) 的上限。}
\end{solutionblock}

\begin{problemblock}
\textbf{例6.} 曲面
\[
(x^2+y^2+z^2)^3=xyz
\]
所围立体的体积为 \(\underline{\hspace{2.5cm}}\).
\end{problemblock}
\begin{solutionblock}
\analysis{用球坐标：
\[
x=\rho\sin\varphi\cos\theta,\quad
y=\rho\sin\varphi\sin\theta,\quad
z=\rho\cos\varphi.
\]
代入曲面方程得
\[
\rho^6=\rho^3\sin^2\varphi\cos\varphi\cos\theta\sin\theta,
\]
即
\[
\rho^3=\sin^2\varphi\cos\varphi\cos\theta\sin\theta.
\]
右端必须为正，因此立体只分布在 \(xyz>0\) 的四个卦限中，且四部分全等。

先算第一卦限。此时
\[
0\le\varphi\le\frac{\pi}{2},\qquad
0\le\theta\le\frac{\pi}{2},
\]
且
\[
0\le\rho\le
\left(\sin^2\varphi\cos\varphi\cos\theta\sin\theta\right)^{1/3}.
\]
第一卦限体积为
\[
V_1=\int_0^{\pi/2}\int_0^{\pi/2}
\int_0^{\rho_{\max}}\rho^2\sin\varphi\,d\rho\,d\theta\,d\varphi
\]
\[
=\frac13\int_0^{\pi/2}\int_0^{\pi/2}
\sin^2\varphi\cos\varphi\cos\theta\sin\theta\cdot\sin\varphi
\,d\theta\,d\varphi
\]
\[
=\frac13\left(\int_0^{\pi/2}\sin^3\varphi\cos\varphi\,d\varphi\right)
\left(\int_0^{\pi/2}\cos\theta\sin\theta\,d\theta\right)
=\frac13\cdot\frac14\cdot\frac12=\frac1{24}.
\]
共有四个全等部分，所以总体积为
\[
V=4V_1=\frac16.
\]}
\answer{\(\displaystyle \frac16\)}
\examnote{方程右端含 \(xyz\) 时，先判断哪些卦限有实体。球坐标下若曲面给出 \(\rho^3=\) 角函数，体积积分会明显简化。}
\end{solutionblock}
